\documentclass[10pt]{beamer}

\usepackage{appendixnumberbeamer}
\usepackage{booktabs}
\usepackage[scale=2]{ccicons}
\usepackage{pgfplots}
\usepackage{xspace}
\usepackage{bookmark}
\usepackage{amssymb}
\usepackage{mathtools}
\usepackage[normalem]{ulem}
\usepackage[T1]{fontenc}
\usepackage[sfdefault,book]{FiraSans}
\usepackage{FiraMono}
\usepackage{fontawesome}
\usepackage{hyperref}
\usepackage{listings}

\renewcommand{\thefootnote}{\fnsymbol{footnote}}
\renewcommand{\MintedPygmentize}{/home/thiago/.local/bin/pygmentize}

\definecolor{light-gray}{gray}{0.95}
\renewcommand\lstlistingname{Código}
\DeclareCaptionFormat{listing} {
  \parbox{\textwidth}{\hspace{-0.2cm}#1#2#3}
}
\DeclareCaptionFont{black}{\color{black}}
\captionsetup[lstlisting]{
  format=listing,
  labelfont=black,
  textfont=black,
  singlelinecheck=true,
  margin=0pt,
  font={tt,footnotesize,bf}
}

\lstset{
  basicstyle=\footnotesize\ttfamily,
  escapeinside={\%*}{*)},
  mathescape=true,
  showspaces=false,
  showtabs=false,
  showstringspaces=false,%
  rulesepcolor=\color{black},
  frame=shadowbox,
  upquote=true,
  literate=
  {á}{{\'a}}1 {é}{{\'e}}1 {í}{{\'i}}1 {ó}{{\'o}}1 {ú}{{\'u}}1
  {Á}{{\'A}}1 {É}{{\'E}}1 {Í}{{\'I}}1 {Ó}{{\'O}}1 {Ú}{{\'U}}1
  {à}{{\`a}}1 {è}{{\`e}}1 {ì}{{\`i}}1 {ò}{{\`o}}1 {ù}{{\`u}}1
  {À}{{\`A}}1 {È}{{\'E}}1 {Ì}{{\`I}}1 {Ò}{{\`O}}1 {Ù}{{\`U}}1
  {ä}{{\"a}}1 {ë}{{\"e}}1 {ï}{{\"i}}1 {ö}{{\"o}}1 {ü}{{\"u}}1
  {Ä}{{\"A}}1 {Ë}{{\"E}}1 {Ï}{{\"I}}1 {Ö}{{\"O}}1 {Ü}{{\"U}}1
  {â}{{\^a}}1 {ê}{{\^e}}1 {î}{{\^i}}1 {ô}{{\^o}}1 {û}{{\^u}}1
  {Â}{{\^A}}1 {Ê}{{\^E}}1 {Î}{{\^I}}1 {Ô}{{\^O}}1 {Û}{{\^U}}1
  {Ã}{{\~A}}1 {ã}{{\~a}}1 {Õ}{{\~O}}1 {õ}{{\~o}}1
  {œ}{{\oe}}1 {Œ}{{\OE}}1 {æ}{{\ae}}1 {Æ}{{\AE}}1 {ß}{{\ss}}1
  {ű}{{\H{u}}}1 {Ű}{{\H{U}}}1 {ő}{{\H{o}}}1 {Ő}{{\H{O}}}1
  {ç}{{\c c}}1 {Ç}{{\c C}}1 {ø}{{\o}}1 {å}{{\r a}}1 {Å}{{\r A}}1
  {€}{{\euro}}1 {£}{{\pounds}}1 {«}{{\guillemotleft}}1
  {»}{{\guillemotright}}1 {ñ}{{\~n}}1 {Ñ}{{\~N}}1 {¿}{{?`}}1
}

\setminted{
  fontsize=\footnotesize,
  style=bw,
  frame=single,
  labelposition=topline,
}

\pretitle{\begin{center}\normalsize\bfseries\MakeUppercase{Programação 2 -- }\MakeUppercase}
\author{\normalsize Prof. Thiago Cavalcante}
\date{\vspace{-4ex}}

\geometry{
  top=0.0cm,
  left=1.0cm,
  right=1.0cm,
  bottom=1.0cm
}

\setlength\columnsep{30pt}

\usetikzlibrary{arrows}

\makeatletter
\AddEnumerateCounter{\PaddingUp}{\two@digits}{A00}
\AddEnumerateCounter{\PaddingDown}{\two@digits}{A00}
\newcommand\PaddingUp[1]{\expandafter\two@digits\csname c@#1\endcsname}
\newcommand\PaddingDown[1]{\PaddingUp{#1}\addtocounter{#1}{-2}}
\makeatother

\makeatletter
\def\enumalphalphcnt#1{\expandafter\@enumalphalphcnt\csname c@#1\endcsname}
\def\@enumalphalphcnt#1{\AlphAlph{#1}}
\makeatother
\AddEnumerateCounter{\enumalphalphcnt}{\@enumalphalphcnt}{aa}

\pagenumbering{gobble}

\title{Programação 2}
\author{Thiago Cavalcante  -- thiago.cavalcante@penedo.ufal.br}
\institute{Universidade Federal de Alagoas -- UFAL \\ Campus Arapiraca \\ Unidade de Ensino de Penedo}
\titlegraphic{\hfill\includegraphics[height=1.5cm]{images/brasao-ufal.eps}}


\subtitle{Aula 5}
\date{27 de novembro de 2019}

\begin{document}

\maketitle

\begin{frame}{Alocação dinâmica}
  \huge
  \textbf{Uma variável é uma \alert{posição na memória previamente* reservada} que pode armazenar um dado}
  \vfill
  *previamente $\rightarrow$ toda variável deve ser \alert{declarada}
\end{frame}

\begin{frame}{Alocação dinâmica}
  \huge
  Nem sempre é possível saber quanta memória um programa vai precisar (ex. funcionários)
\end{frame}

\begin{frame}{Alocação dinâmica}
  \huge
  \textbf{A linguagem C permite \alert{alocar dinamicamente} blocos de memórias usando ponteiros}
  \vfill
  \Large
  \setlength{\leftmargini}{0pt}
  \begin{itemize}
      \item [] Alocar $\rightarrow$ reservar
      \item [] Dinamicamente $\rightarrow$ em tempo de execução
  \end{itemize}
\end{frame}

\begin{frame}{Alocação dinâmica}
  \huge
  O programa requisita, \textbf{em tempo de execução}, um espaço de memória ao computador, o qual retorna um \textbf{endereço para o início} desse espaço (que é armazenado em um \textbf{ponteiro})
\end{frame}

\begin{frame}{Alocação dinâmica}
  \huge
  \textbf{Funções para alocação de memória (\texttt{stdlib.h})}
  \vfill
  \LARGE
  \begin{itemize}
    \item \texttt{sizeof}
    \item \texttt{malloc}
    \item \texttt{calloc}
    \item \texttt{realloc}
    \item \texttt{free}
  \end{itemize}
\end{frame}

\begin{frame}[fragile]{Alocação dinâmica}
  \huge
  \textbf{\texttt{sizeof}}
  \vfill
  Usada para saber quantos bytes ocupa \textbf{um único elemento} de um determinado tipo
  \vfill
  \large
  \begin{lstlisting}[language=C]
printf("%d", sizeof(int)); // resultado: 4
  \end{lstlisting}
\end{frame}

\begin{frame}[fragile]{Alocação dinâmica}
  \huge
  \textbf{\texttt{malloc}}
  \vfill
  \Large
  \begin{lstlisting}[language=C]
void *malloc (unsigned int num);
  \end{lstlisting}
  \vfill
  \large
  \setlength{\leftmargini}{0pt}
  \begin{itemize}
      \item [] Aloca uma quantidade \textbf{\texttt{num}} de bytes de memória e retorna um \textbf{ponteiro genérico (\texttt{void*})} para o início do bloco
      \item [] Programador \textbf{deve atribuir um tipo} ao ponteiro
      \item [] Se houver erro a função retorna \textbf{\texttt{NULL}}
  \end{itemize}
\end{frame}

\begin{frame}[fragile]{Alocação dinâmica}
  \large
  \begin{lstlisting}[language=C]
int i;
int *p;

p = (int *) malloc(5 * sizeof(int));

for(i = 0; i < 5; i++) {
  p[i] = 10 + i;
  // *(p + i) = 10 + i;
}
  \end{lstlisting}
  \vfill
  \large
  \faExclamationTriangle \quad Prestar atenção no tipo de dados e na quantidade
\end{frame}

\begin{frame}[fragile]{Alocação dinâmica}
  \huge
  \textbf{\texttt{calloc}}
  \vfill
  \Large
  \begin{lstlisting}[language=C]
void *calloc (unsigned int num,
              unsigned int size);
  \end{lstlisting}
  \vfill
  \large
  \setlength{\leftmargini}{0pt}
  \begin{itemize}
      \item [] Aloca uma quantidade $\textnormal{\textbf{\texttt{num}}} \times \textnormal{\textbf{\texttt{size}}}$ de bytes de memória e retorna um \textbf{ponteiro genérico (\texttt{void*})} para o início do bloco
      \item [] Inicializa todos os bits do bloco com \textbf{zero}
  \end{itemize}
\end{frame}

\begin{frame}[fragile]{Alocação dinâmica}
  \large
  \begin{lstlisting}[language=C]
int i;
int *p;

// p = (int *) malloc(5 * sizeof(int));
p = (int *) calloc(5, sizeof(int));

for(i = 0; i < 5; i++) {
  p[i] = 10 + i;
  // *(p + i) = 10 + i;
}
  \end{lstlisting}
\end{frame}

\begin{frame}[fragile]{Alocação dinâmica}
  \huge
  \textbf{\texttt{realloc}}
  \vfill
  \Large
  \begin{lstlisting}[language=C]
void *realloc (void *ptr,
               unsigned int num);
  \end{lstlisting}
  \vfill
  \large
  \setlength{\leftmargini}{0pt}
  \begin{itemize}
      \item [] Modifica para \textbf{\texttt{num}} bytes o tamanho da memória previamente alocada e apontada por \textbf{\texttt{ptr}}
  \end{itemize}
\end{frame}

\begin{frame}[fragile]{Alocação dinâmica}
  \large
  \vfill
  \begin{lstlisting}[language=C]
int *p;
p = malloc(5 * sizeof(int));

//Diminui o tamanho do array
p = realloc(p, 3 * sizeof(int));

//Aumenta o tamanho do array
p = realloc(p, 10 * sizeof(int));
  \end{lstlisting}
  \vfill
  Quando o tamanho é maior, o bloco recém-alocado tem valor \textbf{indeterminado} (não ocorre inicialização)
\end{frame}

\begin{frame}[fragile]{Alocação dinâmica}
  \huge
  \textbf{\texttt{free}}
  \vfill
  \Large
  \begin{lstlisting}[language=C]
void free (void *p);
  \end{lstlisting}
  \vfill
  \large
  \setlength{\leftmargini}{0pt}
  \begin{itemize}
      \item [] A memória alocada dinamicamente (\texttt{malloc}, \texttt{calloc} ou \texttt{realloc}) \textbf{não é liberada automaticamente} pelo programa \faExclamationTriangle
      \item [] Essa memória precisa ser liberada com a função \textbf{\texttt{free}}, que recebe o ponteiro para o início do bloco
  \end{itemize}
\end{frame}

\begin{frame}[fragile]{Alocação dinâmica}
  \huge
  \textbf{\alert{Sempre} libere a memória que não for mais utilizar e não deixe ponteiros "soltos"}
  \vfill
  \begin{lstlisting}[language=C]
free(p);
p = NULL;
  \end{lstlisting}
\end{frame}

\begin{frame}[fragile]{Alocação dinâmica}
  \huge
  \textbf{Arrays multidimensionais}
  \vfill
  \LARGE
  \begin{itemize}
    \item Usando array unidimensional
    \item Usando ponteiro para ponteiro
    \item Usando ponteiro para ponteiro + array unidimensional
  \end{itemize}
\end{frame}
\end{document}
